\section{Kustomize}

The kustomize.yaml defines the resources and transformations. \textit{Transformers} transform your manifests by allowing common fields in the resources to be specified in one place.

In case you want multiple variations on a common resource, the files in \texttt{base/} can declare common elements and the files in \texttt{overlays/} will declare differences.

\textbf{Patches} allow you to change (patch) already set values.

Traditionally, when editing ConfigMaps or Secrets, the deployment did not change, therefore no pods were restarted. \textbf{Generators} allow to define values in the \texttt{kustomize.yaml} file, where changes get detected and pods automatically restarted.
\begin{lstlisting}[language=yaml]
# A sample kustomize.yaml file
apiVersion: kustomize.config.k8s.io/v1
kind: Kustomization
resources:
  - deployment.yaml
  - service.yaml
# Transformers:
namePrefix: dev-
namespace: my-namespace
commonLabels:
  app: my-app
# Patch:
patches:
- patch: |-
    - op: replace
      path: /metadata/name
      value: nginx-server
  target:
    kind: Service
    name: nginx-app
# Generator:
configMapGenerator:
  - name: app-config
    literals:
      - LOG_LEVEL=debug
# Adding component from component example below
components:
- ../../components/mysql
\end{lstlisting}

\subsection{Components}

Components encapsulate a set of modifications. Below is an example component which adds a mysql database. This component can then be included in overlays, base directories (see above) or other components.

\begin{lstlisting}[language=yaml]
# Sample component file. Path: components/mysql/kustomization.yaml
apiVersion: kustomize.config.k8s.io/v1alpha1
kind: Component # not Kustomization
resources: # resource files are also under /mysql/
- deployment.yaml
- service.yaml
\end{lstlisting}

\subsection{Commands}

\begin{lstlisting}[language=sh]
kustomize build <dir> # Renders manifests.
kubectl apply -k <dir> # Applies kustomization via kubectl.
\end{lstlisting}

\subsection{Images \& Overlays}
The \textbf{images} transformer changes a container's tag/name without touching the Deployment. \texttt{configMapGenerator}/\texttt{secretGenerator} can also read from \texttt{files:} (not just \texttt{literals:}) to keep config in one place. An overlay pulls in the shared base via \texttt{resources: [../base]}.
\begin{lstlisting}[language=yaml]
images:
- name: myrepo/app   # existing image (no tag)
  newTag: "0.2.0"    # or newName: otherrepo/app
configMapGenerator:
- name: env-cfg
  files: [config.env]
\end{lstlisting}
Render without applying: \texttt{kubectl kustomize <dir>}. Tear down: \texttt{kubectl delete -k <dir>}.

\subsection{Helm vs. Kustomize}
\textbf{Helm:} templating + package manager (variables via \texttt{values.yaml}, releases, repos, rollbacks); good for distributing reusable apps. \\
\textbf{Kustomize:} template-free overlays (patch/merge plain YAML per environment), built into \texttt{kubectl} (\texttt{-k}); good for own manifests across dev/test/prod. The two are often combined (Helm to package, Kustomize to patch).
